% Generated draft: requires mathematical and source review.
\documentclass[11pt]{article}
\usepackage[margin=25mm]{geometry}
\usepackage{fontspec}
\setmainfont{Latin Modern Roman}
\usepackage{amsmath,amssymb,mathtools,unicode-math}
\setmathfont{Latin Modern Math}
\usepackage{enumitem,needspace,xurl,hyperref}
\setlength{\parindent}{0pt}
\setlength{\parskip}{5pt}
\setlength{\emergencystretch}{3em}
\newcommand{\N}{\mathbb{N}}
\newcommand{\Z}{\mathbb{Z}}
\newcommand{\Q}{\mathbb{Q}}
\newcommand{\R}{\mathbb{R}}
\newcommand{\C}{\mathbb{C}}
\newcommand{\F}{\mathbb{F}}
\newcommand{\A}{\mathbb{A}}
\newcommand{\PP}{\mathbb{P}}
\newcommand{\E}{\mathbb{E}}
\newcommand{\eps}{\varepsilon}
\newcommand{\dd}{\,\mathrm{d}}
\begin{document}
\section*{UnsolvedMath ID: 30002805}
\textbf{Source problem code:} catalogue entry
\subsection{Definitions and mathematical statement}
The supplied local registry record for permanent UnsolvedMath ID \(30002805\) is marked \texttt{statement\_status: unrecoverable}. Its verification note says that the extraction combined two illustrative results into a purported open problem. Consequently, the following are incomplete motivating tasks, not a recovered open conjecture or an asserted theorem.

\begin{enumerate}
\item Let \(N\) be the number of nodes, indexed by \(i=1,\ldots,N\). Packets are sent at discrete times over an unreliable communication medium. For node \(i\), let \(p_i\in[0,1]\) denote its stated success probability and let \(q_i\geq0\) denote its required throughput per unit time. The source describes successes as independent, but does not specify the full collection of events to which independence applies. Let \(T\) denote the deadline period; its units and admissible values are not specified.

The motivating task is to characterize exactly the feasible parameter collections
\[
\bigl((p_i,q_i)_{i=1}^{N},T\bigr).
\]
Here feasibility requires a specified packet-arrival model, transmission constraints, scheduling information, deadline convention, and mathematical meaning of throughput. These are absent from the supplied statement. In particular, no choice between expected throughput, almost-sure long-run throughput, or another delivery criterion is justified. Thus the supplied data do not define a unique feasibility predicate.

\item Consider the associated autonomous-traffic system described as communicating unicycle cars on multi-lane roads, including intersections. The motivating task is to design controls guaranteeing no collisions, no road departures, and safe traversal of intersections.

For a fully specified traffic model \(M\), let \(S_M\) denote its set of configurations satisfying its collision and road-containment requirements. The safety component would require a control law under which every trajectory allowed by \(M\), from every initial condition admitted by its assumptions, satisfies
\[
x(t)\in S_M\qquad\text{for every }t\geq0.
\]
This is only a schematic expression of the requested safety guarantee. The record does not specify the unicycle equations, control bounds, vehicle geometry, road and intersection geometry, admissible initial conditions, disturbances, communication assumptions, or boundary-contact conventions. It also does not determine whether safe intersection traversal includes an eventual-passage requirement. These omissions prevent a complete control specification or a precise coupling to the packet problem.
\end{enumerate}

The supplied verification note reports that the official report already provides a sharp packet-throughput feasibility characterization and develops provably safe traffic-control laws. It identifies no separate clean open statement. This source qualification must accompany the motivating tasks above.
\subsection{Short English statement}
The record asks which packet-throughput requirements can be met before deadlines, and how communicating autonomous cars can avoid collisions, remain on the road, and traverse intersections safely. Essential model assumptions are missing. Moreover, the supplied verification note identifies the report's material as illustrative results that were incorrectly combined into an open problem, so no precise unresolved conjecture can be recovered here.
\subsection{Sources}
\begin{itemize}
\item Supplied local registry record, permanent UnsolvedMath ID \(30002805\), title ``Expository Packet-Scheduling and Traffic-Control Results,'' exact source problem code \texttt{catalogue entry}. The record associates the material with 2015 and supplies no individual proposer attribution.
\item \emph{Control Theory: A Mathematical Perspective on Cyber-Physical Systems} (2015), Oberwolfach Reports, pp.\ 597--659, DOI \texttt{10.4171/owr/2015/12}. Supplied source citation and literature-source link: \url{https://doi.org/10.4171/owr/2015/12}. No exact passage locating either formulation is supplied.
\item Mitsch, Ghorbal, Vogelbacher and Platzer, supplied literature-source link: \url{https://arxiv.org/abs/1605.00604}. The record supplies no title or precise theorem reference for this item.
\item The supplied literature assessment was checked on exactly \texttt{2026-08-21}. It labels the record ``partially solved'' and states that the communication subproblem receives a sharp feasibility characterization, that provably safe control laws exist for restricted roadway and intersection models, and that the combined general feasibility and control question remains open. This is supplied dated curation derived from the Oberwolfach material, not an independently verified assessment; its time-sensitive status and cited claims require independent verification.
\item The supplied statement-verification note says that the official report presents two illustrative cyber-physical-systems results rather than unresolved problems, that their motivating questions were joined by extraction, and that no separate clean statement is promoted. The original record is retained publicly for provenance. This qualification prevents treating the supplied open-status assessment as a verified mathematical conclusion.
\end{itemize}
\end{document}
